\section[Desenvolvimento do Projeto]{Desenvolvimento do Projeto}

Esta seção apresenta informações relativas ao desenvolvimento do projeto: localização do código-fonte e documentação, arquitetura e tecnologias utilizadas, e protótipos de tela desenvolvidos. Toda a documentação do projeto pode ser encontrada na wiki do mesmo: \url{https://tools.ages.pucrs.br/projeto-focus/wiki} \cite{wikiFocus}.

\subsection{Repositório do Código Fonte do Projeto}
  Todo o código do projeto encontra-se distribuído em dois repositórios:
\begin{itemize}
    \item código do backend: \url{https://tools.ages.pucrs.br/projeto-focus/backend};
    \item código do frontend: \url{https://tools.ages.pucrs.br/projeto-focus/frontend}.
\end{itemize}

Há, ainda, um terceiro repositório, usado somente para abrigar a wiki do projeto. A wiki encontra-se em: \url{https://tools.ages.pucrs.br/projeto-focus/wiki}.

\subsection{Banco de Dados Utilizado}
  Para a realização do projeto, escolhemos utilizar o banco de dados não relacional Firebase. A razão pela escolha deste modelo de banco de dados se deu principalmente pela escolha do time em aprender novas tecnologias, pela facilidade de integração com o backend da aplicação e também devido ao projeto possuir somente uma entidade a ser armazenada, que é o usuário.

\begin{figure}[H]
    \centering
    \small
    \caption{Modelo Entidade-Relacionamento}
    \fbox{\parbox[c][7cm][c]{0.9\linewidth}{\centering \textit{[Inserir imagem]}}}

    Fonte: Wiki do projeto \cite{wikiFocus}
    \label{fig:focus-mer}
\end{figure}

  Durante a apresentação do projeto, o stakeholder afirmou que não haveria a necessidade de armazenar os dados, pois caso fosse necessário o usuário poderia apenas reescrevê-los. Após uma conversa com o time, conseguimos mostrar como um banco de dados, mesmo que o mais simplificado possível, poderia agregar no seu aplicativo.

\subsection{Arquitetura Utilizada}
  Este projeto foi desenvolvido utilizando a arquitetura MVC (\textit{Model-View-Controller}) e é possível encontrar mais informações na página de arquitetura do projeto: \url{https://tools.ages.pucrs.br/projeto-focus/wiki/-/wikis/arquitetura}.

\subsection{Protótipos das Telas Desenvolvidas}
  Os protótipos deste projeto estão disponíveis no endereço a seguir: \url{https://tools.ages.pucrs.br/projeto-focus/wiki/-/wikis/design}.

\subsection{Tecnologias Utilizadas}
  Nesta seção será abordado um pouco sobre as tecnologias utilizadas ao longo do projeto.

\subsubsection{Flutter}
  O Flutter foi criado pela Google, sendo um framework mais voltado para desenvolvimento de aplicativos móveis. Ele permite a criação de aplicativos compilados nativamente para iOS, Android, Windows, Mac, Linux e Web. Utilizando a linguagem de programação Dart, o Flutter se diferencia de outros frameworks por utilizar uma abordagem de desenvolvimento baseada em \textit{Widgets}, que são elementos gráficos reutilizáveis e personalizáveis \cite{flutter}.

\subsubsection{Firebase}
  O Firebase é uma plataforma de desenvolvimento de aplicativos móveis e web da Google, que fornece diversos recursos para ajudar os desenvolvedores a criar aplicações de forma efetiva, rápida e simples. Ele inclui diversas funcionalidades, incluindo armazenamento de dados em tempo real, autenticação de usuários, hospedagem, mensagens e notificações, análise de dados, entre outros \cite{firebase}.

\subsubsection{Go}
  O Go, também conhecido como Golang, é uma linguagem de programação de código aberto criada pelo Google. Ela foi projetada para oferecer eficiência, simplicidade e escalabilidade, tornando-se uma escolha popular para o desenvolvimento de aplicativos de servidor e sistemas de rede \cite{golang}.
